\subsection{Execution Benchmarks}

\begin{principle} \hlt{Benchmark Matches Trading Need}\\
Each execution benchmark expresses a different risk trade-off: avoid detection, follow market volume, capture price opportunity, or avoid missing liquidity. Benchmark choice should depend on order size, market depth, region, and investment style.
\end{principle}

\begin{taxonomy} \hlt{Regional Benchmark Preference}
\begin{enumerate}[label=\roman*.]
\setlength{\itemsep}{0pt}
\item \hlt{US}: liquidity seeking is natural when HFT is pervasive and stealth matters;
\item \hlt{Asia}: VWAP and PoV help compare local market efficiency and fragmented market conventions;
\item \hlt{Europe}: IS and routing discipline matter because orderbooks and venue interaction are fragmented.
\end{enumerate}
\end{taxonomy}

\begin{table}[htbp]
\centering
\scriptsize
\stepcounter{table}
\renewcommand{\arraystretch}{1.12}
\setlength{\tabcolsep}{4pt}
\begin{tabularx}{\textwidth}{@{}>{\RaggedRight\arraybackslash}p{0.15\textwidth} >{\RaggedRight\arraybackslash}p{0.13\textwidth} >{\RaggedRight\arraybackslash}p{0.25\textwidth} >{\RaggedRight\arraybackslash}p{0.18\textwidth} >{\RaggedRight\arraybackslash}X@{}}
\hline
\textbf{Benchmark} & \textbf{Region} & \textbf{Order characteristics} & \textbf{Market context} & \textbf{Risk hedged} \\
\hline
PoV & Asia & Large order size; more than 10\% of ADV, using average daily consolidated volume & Possible negative news & Rapid news repricing \\
VWAP/TWAP & Asia and Europe & Medium size; 5--15\% of ADV & Unusual volume is negligible & Slow information drift \\
Implementation Shortfall (IS) & Europe and US & Small size; 0--6\% of ADV & Possible price opportunities & Vanishing price opportunity \\
Liquidity Seeker & US & Any size & Stock expected to oscillate around fair value & Liquidity bursts; relative-price dislocations \\
\hline
\end{tabularx}

\vspace{0.25em}
{\footnotesize Table \thetable. Typical uses of benchmarks for trading algorithms.}
\end{table}

\begin{caveat} \hlt{Benchmark Is Not an Objective Alone}\\
The benchmark justifies the comparison, not the whole execution policy. Market conditions, order size, liquidity, and control freedom still determine the usable algorithm.
\end{caveat}
